\chapter{压强：从运动到潜水}\label{ch:g10:pressure}

滑雪者在步行者会陷入的雪面上滑行；一个 \qty{5}{N} 的推力
能让针刺穿拳头压不凹的皮革。
效果却天差地别：缺的那个量正是\emph{\omterm{def:g10:pressure:pressure}{压强}}，
即力除以承受它的面积。本章定义这一量，
随它潜入水池、升入气压计、压进被挤压的注射器——
最终落在海面下四十米处，那里\omterm{def:g10:pressure:pressure}{压强}变得至关重要。

\section{压强：力在面积上的分布}

\begin{definition}[压强]\label{def:g10:pressure:pressure}
当大小为 $F$ 的力（\omterm{def:g10:orders-of-magnitude:unit}{单位} \unit{N}）垂直作用于
面积为 $S$（\omterm{def:g10:orders-of-magnitude:unit}{单位} \unit{m^2}）的表面上时，
该表面受到的\emph{压强}\index{压强}为
\[
P = \frac{F}{S}.
\]
其\omterm{def:g10:orders-of-magnitude:unit}{单位} \unit{N/m^2}，即\emph{帕斯卡}\index{帕斯卡}
（\unit{Pa}）——一个很小的\omterm{def:g10:orders-of-magnitude:unit}{单位}，相当于一个苹果的\omterm{def:g10:universal-gravitation:weight}{重力}
均匀分布在一平方米上，所以实际常用的\omterm{def:g10:orders-of-magnitude:unit}{单位}是 \unit{kPa} 和 \unit{MPa}。
\end{definition}

\begin{example}[高跟鞋与雪鞋]\label{ex:g10:pressure:snowshoe}
一个 \qty{60}{kg} 的人，\omterm{def:g10:universal-gravitation:weight}{重力} $F = 60 \times 9.81 \approx \qty{590}{N}$。
若全压在一只细高跟上，其面积为 \qty{1.0}{cm^2} = \qty{1.0e-4}{m^2}：
$P = 590/\num{1.0e-4} \approx \qty{5.9e6}{Pa} = \qty{5.9}{MPa}$。若穿两只
雪鞋，每只面积为 \qty{0.20}{m^2}：
$P = 590/0.40 \approx \qty{1.5}{kPa}$。\emph{同样}的\omterm{def:g10:universal-gravitation:weight}{重力}，
压在脚跟下的\omterm{def:g10:pressure:pressure}{压强}大了 $4000$ 倍：高跟鞋会在木地板上留下凹坑，
而雪鞋则能让人浮在松雪上。
\end{example}

\begin{omfigure}
\begin{tikzpicture}[scale=1.0, font=\small]
  \draw[thick] (-1.3,0) -- (1.3,0);
  \draw[-Stealth, omThm, very thick] (0,2.1) -- (0,0.6)
    node[midway, right] {$F = \qty{590}{N}$};
  \fill[omDef!30] (-0.09,0) rectangle (0.09,0.55);
  \draw[omDef] (-0.09,0) rectangle (0.09,0.55);
  \node[below] at (0,-0.1) {$S = \qty{1.0}{cm^2}$};
  \node[below] at (0,-0.6) {$P \approx \qty{5.9}{MPa}$};
\end{tikzpicture}\qquad
\begin{tikzpicture}[scale=1.0, font=\small]
  \draw[thick] (-1.7,0) -- (1.7,0);
  \draw[-Stealth, omThm, very thick] (0,2.1) -- (0,0.6)
    node[midway, right] {$F = \qty{590}{N}$};
  \fill[omDef!30] (-1.4,0) rectangle (1.4,0.16);
  \draw[omDef] (-1.4,0) rectangle (1.4,0.16);
  \node[below] at (0,-0.1) {$S = \qty{0.40}{m^2}$};
  \node[below] at (0,-0.6) {$P \approx \qty{1.5}{kPa}$};
\end{tikzpicture}

{\small 同一个力作用于两种面积上：决定表面是否屈服的是\omterm{def:g10:pressure:pressure}{压强}，而非力本身。}
\end{omfigure}

\begin{remark}[集中还是分散]\label{rem:g10:pressure:tools}
一根缝衣针，以 \qty{5}{N} 的力压在面积为 \qty{0.01}{mm^2} =
\qty{1e-8}{m^2}，产生 $\qty{5e8}{Pa} = \qty{500}{MPa}$ 的\omterm{def:g10:pressure:pressure}{压强}——
足以撑开皮革纤维。刺穿类工具集中力；承重
设计（滑雪板、履带、地基）则分散力。体育运动
其实是对\omterm{def:g10:pressure:pressure}{压强}的巧妙运用：冰刀咬入冰面，滑雪板则滑行于雪上。
\end{remark}

\section{静止液体中的压强}

为何水在越深处压得越厉害？隔离出一段假想的竖直
水柱，其横截面积为 $S$，自液面延伸至深度 $h$：
体积为 $Sh$，质量为 $\rho S h$（$\rho$ 为液体密度，\omterm{def:g10:orders-of-magnitude:unit}{单位}
为 \unit{kg/m^3}），\omterm{def:g10:universal-gravitation:weight}{重力}为 $\rho S h g$。柱底支撑着
该\omterm{def:g10:universal-gravitation:weight}{重力}，\emph{加上}大气作用于顶面的推力 $P_0 S$。
将总作用力除以 $S$，可知深度 $h$ 处的\omterm{def:g10:pressure:pressure}{压强}应为
$P_0 + \rho g h$——事实确实如此，液体中处处成立。

\begin{theorem}[深度处的压强]\label{thm:g10:pressure:depth}
设密度为 $\rho$ 的液体处于静止状态且与大气相通，则
液面以下深度 $h$ 处的\omterm{def:g10:pressure:pressure}{压强}为
\[
P = P_0 + \rho g h,
\]
其中 $P_0$ 为液面处的\omterm{def:g10:pressure:atmospheric}{大气压强}。该\omterm{def:g10:pressure:pressure}{压强}
仅取决于深度，与容器的形状无关；同一点处，
液体向各个方向的\omterm{def:g10:pressure:pressure}{压强}都相等。
\end{theorem}
\admitted

\begin{remark}
水柱论证只能使这一定律显得可信；
严格的推导（适用于任意形状及各个方向）见本科一年级分册。水自身产生的
这一项 $\rho g h$ 即\emph{表压}\index{表压}，也就是潜水员
压力表所显示的读数。
\end{remark}

\begin{example}[十米水深]\label{ex:g10:pressure:tenmetres}
在淡水中（$\rho = \qty{1000}{kg/m^3}$），当 $h = \qty{10}{m}$ 时：
$\rho g h = 1000 \times 9.81 \times 10 \approx \qty{98}{kPa}$，因此
$P = 101 + 98 \approx \qty{199}{kPa}$。十米水深几乎恰好增加
一个大气压：这是潜水员的经验法则，即\emph{每
下潜十米，\omterm{def:g10:pressure:pressure}{压强}多一个大气压}。
\end{example}

\begin{omfigure}
\begin{tikzpicture}[scale=1.0, font=\small]
  \shade[top color=omDef!8, bottom color=omDef!50] (0,0) rectangle (3.4,-3.4);
  \node[above] at (1.7,0.05) {液面：$P_0 = \qty{101}{kPa}$};
  \draw[-Stealth, omThm, thick] (1.7,-1.1) -- ++(1.0,0) node[right] {\qty{199}{kPa}};
  \draw[-Stealth, omThm, thick] (1.7,-2.2) -- ++(1.5,0) node[right] {\qty{297}{kPa}};
  \draw[-Stealth, omThm, thick] (1.7,-3.3) -- ++(2.0,0) node[right] {\qty{395}{kPa}};
  \draw (-0.15,-1.1) -- (0,-1.1) node[left=8pt] {\qty{10}{m}};
  \draw (-0.15,-2.2) -- (0,-2.2) node[left=8pt] {\qty{20}{m}};
  \draw (-0.15,-3.3) -- (0,-3.3) node[left=8pt] {\qty{30}{m}};
\end{tikzpicture}

{\small 静止淡水中，\omterm{def:g10:pressure:pressure}{压强}随深度线性增加，每深 10 米增加一个大气压。}
\end{omfigure}

\section{空气的海洋}

\begin{definition}[大气压强]\label{def:g10:pressure:atmospheric}
我们生活在空气海洋的底部，空气的重量产生了
\emph{大气压强}\index{大气压强}，海平面处为
\[
P_0 \approx \qty{101}{kPa} = \qty{1.01e5}{Pa}.
\]
每平方厘米上的压力为 $F = \num{1.01e5} \times \num{1e-4}
\approx \qty{10}{N}$，相当于一\omterm{def:g10:orders-of-magnitude:unit}{千克}质量的物体的重量作用于
你皮肤、桌面及万物的每 \unit{cm^2} 上。
\end{definition}

\begin{example}[看不见的负荷]\label{ex:g10:pressure:sheet}
一张 A4 纸（\qty{21.0}{cm} $\times$ \qty{29.7}{cm}，故
$S \approx \qty{0.0624}{m^2}$）受到上方空气施加的力
$F = \num{1.01e5} \times 0.0624 \approx \qty{6.3e3}{N}$——相当于
一辆小汽车的重量。纸张不会被撕破：下方的空气
\emph{向上}推得同样用力。只有当一侧失去大气时，我们才会
注意到大气的存在——例如吸盘、吸管、飞机客舱。
\end{example}

\begin{remark}[托里拆利气压计]\label{rem:g10:pressure:torricelli}
1643 年，托里拆利将一根封闭的水银管倒插于
水银槽中，水银柱降至使 $\rho g h = P_0$ 成立的高度：
$h = \num{1.01e5}/(13600 \times 9.81) \approx \qty{0.76}{m}$，其上方
留下了人类历史上第一个真空。水银柱的高度随天气变化，
并随海拔升高而减小：气压计称量的是你上方的空气（若改用水，
则需 \qty{10.3}{m}：见\cref{exo:g10:pressure:9}）。
\end{remark}

\section{气体：碰撞与玻意耳定律}

气体同样会压迫其容器：它是由大量分子组成的
做不停息的无规则运动，每次碰撞都给器壁一个微小的向外
推力，每秒数以亿亿计的撞击汇成了
我们测得的稳定作用力，即气压。将气体压缩至
原体积的一半：分子撞击的频率加倍，因此\omterm{def:g10:pressure:pressure}{压强}应当
加倍。事实确实如此。

\begin{proposition}[玻意耳定律]\label{prop:g10:pressure:boyle}
\index{玻意耳定律}
一定质量的气体，在温度保持不变时，\omterm{def:g10:pressure:pressure}{压强}和
体积成反比：
\[
P \, V = \text{常量},
\qquad\text{即}\qquad
P_1 V_1 = P_2 V_2 .
\]
\end{proposition}
\admitted

\begin{remark}
这是本阶段的一个实验事实——本科一年级分册
由碰撞图像推导出该结论。实验装置（密封的注射器内充有空气、
\omterm{def:g10:pressure:pressure}{压强}传感器，\qty{20}{\degreeCelsius}）：
\begin{center}
\begin{tabular}{lccccc}
$V$ (\unit{mL}) & 60 & 50 & 40 & 30 & 20 \\
$P$ (\unit{kPa}) & 100 & 120 & 150 & 200 & 300 \\
$P \times V$ (\unit{kPa.mL}) & 6000 & 6000 & 6000 & 6000 & 6000
\end{tabular}
\end{center}
\end{remark}

\begin{omfigure}
\begin{tikzpicture}
\begin{axis}[omaxis, width=7.6cm, height=5.6cm,
    xlabel={$V$ (\unit{mL})}, ylabel={$P$ (\unit{kPa})},
    xmin=0, xmax=70, ymin=0, ymax=340,
    domain=18.5:68, samples=100]
  \addplot[omThm, thick] {6000/x};
  \addplot[only marks, mark=*, omDef]
    coordinates {(60,100) (50,120) (40,150) (30,200) (20,300)};
\end{axis}
\end{tikzpicture}

{\small 注射器数据落在双曲线 $P = 6000/V$ 上：体积减半，\omterm{def:g10:pressure:pressure}{压强}加倍。}
\end{omfigure}

\begin{example}[压缩的注射器]\label{ex:g10:pressure:boyle}
空气在 \qty{100}{kPa} 下占 \qty{4.0}{L}；缓慢压缩（使
温度保持不变）至 \qty{250}{kPa}，其体积为
$V_2 = P_1 V_1 / P_2 = 400/250 = \qty{1.6}{L}$.
\end{example}

\section{潜水：压强的应用}

\begin{method}[深度处的绝对压强]\label{met:g10:pressure:dive}
潜水员在海中深度 $h$ 处（$\rho \approx \qty{1025}{kg/m^3}$，
故每 \qty{10}{m} 深度 $\rho g h \approx \qty{101}{kPa}$），计算
$P = P_0 + \rho g h$——在玻意耳定律中，始终使用这一
\emph{绝对压强}\index{绝对压强}，切勿仅用\omterm{thm:g10:pressure:depth}{表压}项。潜水员
在 \qty{10}{m} 处承受 $2$ 个大气压，在 \qty{20}{m} 处承受 $3$ 个大气压，在
\qty{40}{m} 处承受 $5$ 个大气压。
\end{method}

\begin{omfigure}
\begin{tikzpicture}[scale=1.0, font=\small]
  \shade[top color=omDef!8, bottom color=omDef!55] (0,0) rectangle (2.0,-4.4);
  \draw[thick] (2.0,0) -- (2.0,-4.4);
  \draw[thick, omDef] (-0.4,0) -- (2.4,0);
  \draw (1.85,0) -- (2.15,0) node[right] {\qty{0}{m}：\qty{101}{kPa}（$1$ 个大气压）};
  \draw (1.85,-1.1) -- (2.15,-1.1) node[right] {\qty{10}{m}：\qty{202}{kPa}（$2$ 个大气压）};
  \draw (1.85,-2.2) -- (2.15,-2.2) node[right] {\qty{20}{m}：\qty{302}{kPa}（$3$ 个大气压）};
  \draw (1.85,-3.3) -- (2.15,-3.3) node[right] {\qty{30}{m}：\qty{403}{kPa}（$4$ 个大气压）};
  \draw (1.85,-4.4) -- (2.15,-4.4) node[right] {\qty{40}{m}：\qty{503}{kPa}（$5$ 个大气压）};
  \fill[omThm] (0.9,-2.15) circle (0.09);
  \draw[omThm, thick] (0.9,-2.24) -- (0.9,-2.6) (0.9,-2.35) -- (0.7,-2.5)
    (0.9,-2.35) -- (1.1,-2.5) (0.9,-2.6) -- (0.78,-2.85) (0.9,-2.6) -- (1.02,-2.85);
\end{tikzpicture}

{\small 海水中的\omterm{met:g10:pressure:dive}{绝对压强}：水面处为一个大气压，每下潜 10 米再增加一个大气压。}
\end{omfigure}

\begin{remark}[切勿屏住呼吸]\label{rem:g10:pressure:breath}
减压器以环境\omterm{def:g10:pressure:pressure}{压强}供气，肺部在任何深度都能正常充气，
但若在 \qty{20}{m} 深处（$\qty{302}{kPa}$）充入 \qty{6.0}{L} 的空气，
屏住呼吸后浮出水面，玻意耳定律要求
$6.0 \times 302/101 \approx \qty{18}{L}$——约为胸腔容量的三倍，
肺部早在此之前就会破裂，即使只上升几米也是如此。
因此水肺潜水的第一条规则是：\emph{上升时绝不可屏住呼吸}
——保持呼吸，膨胀的空气便会自行排出。自由潜水员
是安全的：其所吸入的空气取自水面，因此上升时
空气只会重新膨胀至原有体积。
\end{remark}

\begin{remark}[减压病]\label{rem:g10:pressure:deco}
呼吸高压空气还会使血液中溶解额外的氮气，
如同密封汽水瓶中的气体；上升过快时，氮气便会在体内形成
气泡——即减压病。因此上升宜缓慢
并在浅水处停留（每分钟约上升 \qty{10}{m}），
使气体经肺部缓慢排出。其理论留待本科一年级分册讲解，
而行为准则则属于每一位潜水员。
\end{remark}

\section{练习}

\begin{exercise}[$\star$]\label{exo:g10:pressure:1}
一只重 \qty{600}{N} 的木箱先后平放在
$\qty{0.50}{m} \times \qty{0.30}{m}$ 的面上，再放在
$\qty{0.30}{m} \times \qty{0.20}{m}$ 的面上。求木箱对地面的两个\omterm{def:g10:pressure:pressure}{压强}。
哪种放置方式在松软的沙地上更好？
\end{exercise}

\begin{exercise}[$\star$]\label{exo:g10:pressure:2}
将 \qty{0.25}{MPa} 分别换算为 \unit{kPa} 和 \unit{Pa}。下列两种情形中，
哪一种\omterm{def:g10:pressure:pressure}{压强}更大：\qty{3.0}{N} 作用于 \qty{2.0}{cm^2}，
还是 \qty{60}{N} 作用于 \qty{500}{cm^2}？
\end{exercise}

\begin{exercise}[$\star$]\label{exo:g10:pressure:3}
求游泳池池底处的\omterm{thm:g10:pressure:depth}{表压} $\rho g h$ 和
\omterm{met:g10:pressure:dive}{绝对压强}。游泳池的水深为 \qty{3.0}{m}，
（$\rho = \qty{1000}{kg/m^3}$，$P_0 = \qty{101}{kPa}$）。
\end{exercise}

\begin{exercise}[$\star$]\label{exo:g10:pressure:4}
空气在 \qty{100}{kPa} 下占体积 \qty{4.0}{L}，温度保持不变。
求其在 \qty{250}{kPa} 时的体积？在 \qty{50}{kPa} 时呢？
\end{exercise}

\begin{exercise}[$\star$]\label{exo:g10:pressure:5}
一名质量为 \qty{65}{kg} 的滑冰者单脚滑行，冰刀接触面积为
\qty{12}{cm^2}；随后双脚站立，鞋底总面积为 \qty{360}{cm^2}。
求两种情形下的\omterm{def:g10:pressure:pressure}{压强}及其比值。
\end{exercise}

\begin{exercise}[$\star\star$]\label{exo:g10:pressure:6}
用拇指以 \qty{20}{N} 的力按压图钉。钉帽面积为
\qty{1.2}{cm^2}，针尖面积为 \qty{0.010}{mm^2}。求拇指与墙面受到的\omterm{def:g10:pressure:pressure}{压强}，
并说明为何只有墙面被扎穿？
\end{exercise}

\begin{exercise}[$\star\star$]\label{exo:g10:pressure:7}
一艘潜艇的舱门面积为 \qty{0.50}{m^2}，位于水下 \qty{30}{m} 深处的
海水中（$\rho = \qty{1025}{kg/m^3}$）；舱内空气处于
\omterm{def:g10:pressure:atmospheric}{大气压强}下。求使舱门保持关闭的合力，以及
\omterm{def:g10:universal-gravitation:weight}{重力}恰好与之相等的质量。
\end{exercise}

\begin{exercise}[$\star\star$]\label{exo:g10:pressure:8}
一个传感器位于未知液体中 \qty{5.00}{m} 深处，
测得\omterm{thm:g10:pressure:depth}{表压}为 \qty{66.2}{kPa}。求该液体的密度，并判断其是否为淡水？
\end{exercise}

\begin{exercise}[$\star\star$]\label{exo:g10:pressure:9}
若用水（$\rho = \qty{1000}{kg/m^3}$）代替水银
（$\rho = \qty{13600}{kg/m^3}$）制作托里拆利气压计，
其高度应为多少？由 $\rho g h = P_0$ 分别求出两个高度；
为何水银更为合适？
\end{exercise}

\begin{exercise}[$\star\star$]\label{exo:g10:pressure:10}
利用 \cref{prop:g10:pressure:boyle} 中的注射器数据
（$PV = \qty{6000}{kPa.mL}$），预测\omterm{def:g10:pressure:pressure}{压强}为 \qty{240}{kPa} 时的体积，以及
体积为 \qty{75}{mL} 时的\omterm{def:g10:pressure:pressure}{压强}。$P$ 对 $V$ 的图像是什么曲线？
$P$ 对 $1/V$ 的图像又是什么？
\end{exercise}

\begin{exercise}[$\star\star$]\label{exo:g10:pressure:11}
一名潜水员在 \qty{20}{m} 深处释放了一个体积为 \qty{0.50}{cm^3} 的气泡。
求气泡刚接近水面时的体积？（温度保持不变；海水
每 \qty{10}{m} 深度增加 \qty{101}{kPa} \omterm{def:g10:pressure:pressure}{压强}。）
\end{exercise}

\begin{exercise}[$\star\star\star$]\label{exo:g10:pressure:12}
一个面积为 \qty{12}{cm^2} 的吸盘被平压在天花板上。
\begin{enumerate}
  \item 若吸盘内为完全真空，大气对吸盘的作用力为多大，
        并可悬挂多大质量的物体？
  \item 实际吸盘内仍有\omterm{def:g10:pressure:pressure}{压强}为 \qty{20}{kPa} 的空气，此时可悬挂多大质量？
\end{enumerate}
\end{exercise}

\begin{exercise}[$\star\star\star$]\label{exo:g10:pressure:13}
一只水肺气瓶内装有 \qty{12}{L}、\omterm{def:g10:pressure:pressure}{压强}为 \qty{200}{bar} 的空气，
（\qty{1}{bar} = \qty{100}{kPa} $\approx 1$ 个大气压）；\qty{50}{bar} 的部分
须留作备用。潜水员在环境\omterm{def:g10:pressure:pressure}{压强}下每分钟呼吸
\qty{15}{L} 空气。可用的水面\omterm{def:g10:pressure:pressure}{压强}下空气有多少升？
气瓶在 \qty{30}{m} 深处（\omterm{met:g10:pressure:dive}{绝对压强}
$\approx \qty{4}{bar}$）能持续多长时间？在 \qty{10}{m} 深处呢？
\end{exercise}

\begin{exercise}[$\star\star\star$]\label{exo:g10:pressure:14}
自由潜水员的肺在水面时容纳 \qty{6.0}{L} 气体，
且不能被压缩到低于残余体积 \qty{1.5}{L}。问：海水深度为多少时
（$\rho = \qty{1025}{kg/m^3}$）会达到这一极限？（此极限曾被视为
自由潜水的极限；如今纪录已超过 \qty{100}{m}，血液
转移进入胸腔，以补上缺失的体积。）
\end{exercise}

\begin{exercise}[$\star\star\star$]\label{exo:g10:pressure:15}
在液压升降机中，液体将\omterm{def:g10:pressure:pressure}{压强}不变地从一个
小活塞（$s = \qty{2.0}{cm^2}$）传递到一个
大活塞（$S = \qty{400}{cm^2}$）上，该大活塞承载一辆 \qty{1200}{kg} 的汽车。
\begin{enumerate}
  \item 小活塞上需要多大的力才能托住汽车？
        此时液体承受的\omterm{def:g10:pressure:pressure}{压强}是多少？
  \item 要使汽车升高 \qty{2.0}{cm}，小活塞
        须移动多远？并比较两个活塞所做的功。
\end{enumerate}
\end{exercise}

\section{问题：一口气下潜，一口气上升}

\begin{problem}[{周末问题——一口气下潜，一口气上升：
自由潜水与水肺上升的物理，从面镜压迫到
“最后十米最危险”这一规则}]
\label{pb:g10:pressure:1}
莉娜进行自由潜水：在水面吸一口气，下潜至 \qty{30}{m} 再返回。
马可在她身边进行水肺潜水，用气瓶呼吸。
两人处于同一片海域（$\rho = \qty{1025}{kg/m^3}$，$g = \qty{9.81}{N/kg}$，
$P_0 = \qty{101}{kPa}$），肺容量同为 \qty{6.0}{L}，危险却截然相反——
二者均由 $P = P_0 + \rho g h$ 及 $PV$ 为常量（温度
全程保持不变）。

\medskip
\noindent\textbf{第一部分——水的重量。}
\begin{enumerate}
  \item 计算海水的 $\rho g$，\omterm{def:g10:orders-of-magnitude:unit}{单位}为 \unit{kPa} 每米，并
        检验经验法则：每 \qty{10}{m} 约相当于一个大气压。
  \item 计算在 \qty{10}{m}、\qty{20}{m}、
        \qty{30}{m} 和 \qty{40}{m} 深度处的\omterm{met:g10:pressure:dive}{绝对压强}。
  \item 将各结果表示为水面\omterm{def:g10:pressure:pressure}{压强}的倍数。
  \item 莉娜的面镜面积为 \qty{150}{cm^2}。在 \qty{30}{m} 深处，若面镜内空气
        仍保持水面\omterm{def:g10:pressure:pressure}{压强}，则水施加的合力
        将为多大？（潜水员经鼻子向面镜内呼气，
        以防止这种\emph{面镜压迫}。）
  \item 鼓膜面积约为 \qty{0.60}{cm^2}。若中耳保持水面\omterm{def:g10:pressure:pressure}{压强}，计算
        在 \qty{3.0}{m} 深处鼓膜所受的合力；
        为何在普通泳池中耳朵也会疼，
        “平衡耳压”应达到什么目的？
\end{enumerate}

\medskip
\noindent\textbf{第二部分——一口气下潜。}
莉娜从水面出发，肺内有 \qty{6.0}{L} 空气。
\begin{enumerate}[resume]
  \item 用玻意耳定律，以\omterm{met:g10:pressure:dive}{绝对压强}计算她
        在 \qty{10}{m}、\qty{20}{m} 和 \qty{30}{m} 深处的肺容量。
  \item 她的肺容量减半时位于多深处？
  \item 肺容量是否与深度成正比？描述
        $V$ 随\omterm{met:g10:pressure:dive}{绝对压强} $P$ 变化的曲线。
  \item 她的残余体积为 \qty{1.5}{L}：肺
        不能再缩小。证明该极限在 \qty{30}{m} 附近达到。
  \item 自由潜水纪录仍超过 \qty{100}{m}：什么
        补上了缺失的体积？（还有什么能流入胸腔？）
\end{enumerate}

\medskip
\noindent\textbf{第三部分——一口气上升。}
在 \qty{30}{m} 深处，马可的调节器为他 \qty{6.0}{L} 的肺充气。
\begin{enumerate}[resume]
  \item 调节器输送这些空气时的\omterm{def:g10:pressure:pressure}{压强}是多少？这些空气
        的密度是水面空气的多少倍？
  \item 马可惊慌失措，屏住呼吸上升至 \qty{10}{m} 深处：
        他肺内被困的空气在该处将占据多大体积？
  \item 到达水面时体积为多少——是其肺容量的多少倍？
  \item 陈述水肺潜水员的第一条规则，以及
        为何正常呼吸即可消除危险。
  \item 莉娜也从 \qty{30}{m} 深处屏住呼吸上升至水面，
        她的肺却完全安全。请解释这种不对称性。
\end{enumerate}

\medskip
\noindent\textbf{第四部分——气泡与最后十米。}
\begin{enumerate}[resume]
  \item 马可在 \qty{40}{m} 深处释放一个 \qty{1.0}{cm^3} 的气泡。
        计算其在 \qty{30}{m}、\qty{20}{m}、\qty{10}{m}
        及紧邻水面下方处的体积。
  \item 对上升过程中每个 \qty{10}{m} 阶段，计算气泡增大的
        倍数。气泡在哪一段增大得最快？
  \item 潜水员上升速度至多为每分钟 \qty{10}{m}，从
        \qty{40}{m} 上升需要多长时间？在哪一分钟内，任何被困或溶解的
        气体膨胀最为显著？
  \item 在深水中，马可的血液溶解了额外的氮气，犹如
        封口汽水瓶中的气体。过快上升会造成什么，以及
        慢速上升与 \qty{5}{m} 安全停留如何防止这种情况？
  \item 总结：每人用一句话给出自由潜水员的结论、
        水肺潜水员的结论，以及定量的理由：
        为何任何上升的最后十米最值得重视。
\end{enumerate}
\end{problem}
